\chapter{Experiments}

\section{Evaluation Dataset}

\begin{table*}[t]
\centering
%\setlength{\tabcolsep}{14pt}
%\renewcommand{\arraystretch}{1.15}
\caption{Comparison of DataBench with existing OTQA datasets.
DataBench contains substantially larger individual tables than existing OTQA benchmarks.}
\label{tab:dataset_comparison}
\begin{tabular}{lrr}
\toprule
\textbf{Dataset} & 
%\textbf{\# Table} &
\textbf{Avg. \# Rows} &
\textbf{Avg. \# Cols.} \\
\midrule
%Spider~\cite{yu-etal-2018-spider} & 5.1 & - & 27.6\\
%\emph{Open-Domain Table QA}\\
% TableBench~\cite{wu2025tablebench} 
%& 3,681
% & 7 & 17\\
NQ-TABLES \cite{nqtables} %& 
    & 11 & 6 \\
Open-WikiTable \cite{openwikitable} %& 
    & 14 & 6 \\
% DataBench \cite{databench} & \\
\midrule
Derived DataBench~\cite{databench} &  \textbf{46,325} & \textbf{23} \\
\bottomrule
\end{tabular}
\end{table*}
\begin{table}[t]
\centering
\caption{Distribution of serialized table sizes in the derived DataBench corpus.}
\label{tab:databench_table_size}
\begin{tabular}{lr}
\toprule
\textbf{Statistic} & \textbf{Size (Million Tokens)} \\
\midrule
Mean            & 6.40 \\
Median          & 0.76 \\
75th percentile & 5.05 \\
90th percentile & 13.46 \\
95th percentile & 22.87 \\
Maximum         & 125.04 \\
\bottomrule
\end{tabular}
\end{table}

The lack of an OTQA benchmark with substantially large individual tables motivates us to derive an open-domain evaluation setting from an existing large-table CTQA benchmark.
Specifically, we use DataBench~\cite{databench}, which was originally designed for table question answering over real-world tabular datasets.

Table~\ref{tab:dataset_comparison} compares the average table size of DataBench with existing OTQA benchmarks.
NQ-TABLES~\cite{nqtables} and Open-WikiTable~\cite{openwikitable} contain relatively small individual tables, with approximately $11$ and $14$ rows per table on average, respectively.
In contrast, DataBench contains approximately $46$K rows and $23$ columns per table on average, providing a substantially larger table scale for evaluating OLTQA.

We construct the open-domain corpus using all $80$ tables available from the SemEval-2025 Task 8~\cite{semeval2025task8} release, and evaluate end-to-end OLTQA performance on the $522$ questions from the test set.
The original DataBench consists of $65$ real-world tables, while SemEval-2025 Task 8 additionally provides a test set containing $15$ tables and $522$ query--answer pairs.
Unlike the original CTQA setting, where each question is associated with a given gold table, our open-domain setting pools all $80$ tables into a shared corpus and requires the system to retrieve tables before producing the final answer.

We further analyze the serialized table sizes to determine whether the DataBench corpus satisfies the large-table definition introduced in Section~\ref{sec:large_table}.
Each table is serialized as CSV with column headers and without row indices, and token counts are computed using the \texttt{o200k\_base} BPE tokenizer~\footnote{\url{https://github.com/openai/tiktoken}}.
As shown in Table~\ref{tab:databench_table_size}, $47.5\%$ of the tables contain at least one million tokens.
The mean serialized table size is approximately $6.4$ million tokens.
These results show that a substantial proportion of the corpus consists of large tables under our operational definition, making the derived DataBench suitable for evaluating OLTQA.

\paragraph{Open-Domain Retrievability Check.}
Because the DataBench questions were originally written under a closed-domain setting, pooling all tables into a shared corpus may introduce ambiguity for questions that do not contain sufficient information to distinguish their gold tables.
For example, a question such as ``What is the average age?'' may be answerable when the target table is given, but may not uniquely identify the intended table when multiple tables contain an age attribute.
We therefore conduct a retrievability check using our hybrid retriever to examine whether the annotated gold tables remain practically retrievable in the derived open-domain setting.
On the $522$ test questions, the retriever achieves Recall@1, Recall@5, and Recall@10 of $0.65$, $0.83$, and $0.91$, respectively.
In particular, the annotated gold table can be recovered within the top-$10$ candidates for $91\%$ of the test questions, indicating that the gold tables remain practically retrievable for most questions in the derived open-domain setting.
Nevertheless, retrieval success does not rule out ambiguity among multiple plausible tables; we further examine cases of ambiguity in the error analysis in Section~\ref{sec:error_analysis}.

\section{Evaluation Metrics}

We evaluate end-to-end OLTQA performance using answer accuracy and table retrieval performance using Recall@$k$.

For end-to-end evaluation, we follow the official DataBench evaluator used in SemEval-2025 Task 8~\cite{semeval2025task8} and report answer accuracy as the primary evaluation metric.
Rather than relying on generic exact string matching, the evaluator applies type-specific comparison rules according to the expected answer type.
Boolean answers are normalized to equivalent true/false expressions, categorical answers are compared after normalization, and numerical answers are normalized and truncated to two decimal places before comparison.
For list-valued answers, the evaluator applies the corresponding element-level normalization and compares the lists without considering element order.
The final accuracy is computed as the proportion of questions whose predicted answers are judged correct by the evaluator.

% For table retrieval, we report Recall@$k$, which measures the proportion of questions for which the annotated gold table appears among the top-$k$ retrieved candidate tables.

\section{Experimental Setup}
\subsection{LLM Configuration}

We evaluate the OLTQA pipelines across four LLM backebones to examine whether the observed performance trends generalize across different model families and scales.
Specifically, we evaluate GPT-5.6-Sol~\footnote{\url{https://developers.openai.com/api/docs/models/gpt-5.6-sol}},
Gemini-3.7-Flash~\footnote{\url{https://deepmind.google/models/model-cards/gemini-3-7-flash/}},
Qwen3.5-9B~\footnote{\url{https://huggingface.co/Qwen/Qwen3.5-9B}},
and Qwen3.5-4B~\footnote{\url{https://huggingface.co/Qwen/Qwen3.5-4B}}.

For models that expose configurable reasoning effort, we use the minimum available setting to ensure a consistent comparison across OLTQA pipelines, rather than optimizing each pipeline for its best possible performance.
Specifically, GPT-5.6-Sol uses \texttt{reasoning\_effort=none}, while Gemini-3.7-Flash uses \texttt{thinking=low}.
The same LLM configuration is used across all LLM-based components within each OLTQA pipeline.

The Qwen models are served locally through vLLM with a maximum context length of $32$K tokens due to GPU memory constraints.
For API-based models, we do not impose additional limitations on context length beyond those of the respective APIs.

\subsection{Retriever Configuration}

Our table retriever combines BM25 sparse retrieval with dense retrieval using \texttt{sentence-transformers/all-mpnet-base-v2}~\footnote{\url{https://huggingface.co/sentence-transformers/all-mpnet-base-v2}}.
We use the same dense embedding model employed in one stage of CRAFT's cascaded retrieval pipeline to reduce differences attributable to the choice of dense embedding model when comparing the retrieval methods.

For the baseline methods, we retain their original retrieval configurations.
OpenTab uses BM25 for initial table retrieval.
CRAFT employs its cascaded retrieval pipeline, which progressively filters candidate tables using SPLADE~\cite{splade}, \texttt{sentence-transformers/all-mpnet-base-v2}, and \texttt{text-embedding-3-large}~\footnote{\url{https://developers.openai.com/api/docs/models/text-embedding-3-large}} across successive retrieval stages.

\section{Baselines}

We compare our method with two representative OTQA pipelines, a retrieval-enhanced variant, and a closed-domain setting.

\textit{OpenTab}~\cite{opentab} is an OTQA pipeline that combines table retrieval with SQL-guided reranking and reasoning.
We include OpenTab because it explicitly targets scalability to large tables in OTQA, making it a particularly relevant baseline for our OLTQA setting.

\textit{OpenTab+r} is a variant of \textit{OpenTab}.
We replace OpenTab's initial retriever with our table retriever described in Section~\ref{sec:method}.
All subsequent OpenTab components remain unchanged.
This controlled variant allows us to examine the effect of improving the initial table retrieval stage while keeping the remaining OpenTab pipeline unchanged.

\textit{CRAFT}~\cite{craft} is a recent OTQA method that emphasizes high-quality table retrieval through a cascaded retrieval pipeline before providing a compact table context to an LLM for answer generation.
We include CRAFT because its retrieval-oriented design and multi-candidate reasoning pipeline provide a strong comparison.

\textit{Closed-domain setting} directly provides the gold table to the original TableRAG~\cite{tablerag} reasoner, removing corpus-level table retrieval errors.
This setting serves as a diagnostic reference for evaluating reasoning stage performance when the gold table is given.

Implementation details for the evaluated methods are provided in Section~\ref{sec:implementation_detail}.

\section{Candidate Size Selection}
\begin{table}[t]
\centering
\caption{Development-set retrieval performance for candidate-size selection. Marginal gain measures the increase in Recall@$k$ per additional candidate table between consecutive reported candidate sizes.}
\label{tab:candidate_size_selection}
\begin{tabular}{ccccc}
\toprule
& \multicolumn{2}{c}{\textbf{Ours / OpenTab+r}} 
& \multicolumn{2}{c}{\textbf{CRAFT}} \\
\cmidrule(lr){2-3}
\cmidrule(lr){4-5}
\textbf{$k$}
& \textbf{Recall@$k$}
& \textbf{Marginal Gain}
& \textbf{Recall@$k$}
& \textbf{Marginal Gain} \\
\midrule
1  & 0.61 & --     & 0.55 & --     \\
3  & 0.85 & 0.1203 & 0.76 & 0.1050 \\
5  & 0.92 & 0.0313 & 0.83 & 0.0350 \\
10 & 0.97 & 0.0100 & 0.91 & 0.0160 \\
20 & 0.98 & 0.0016 & 0.96 & 0.0050 \\
\bottomrule
\end{tabular}
\end{table}
The role of the candidate size $k$ differs across methods because of their different retrieval and reasoning architectures.
For our method and CRAFT, $k$ denotes the number of candidate tables retained for reasoning.
In OpenTab and OpenTab+r, $k$ instead denotes the number of initially retrieved candidates passed to the reranking stage, after which only the top-ranked table is used for reasoning.

We select $k$ based on the saturation of gold-table coverage on the DataBench development set rather than end-to-end answer accuracy.
The development set contains $320$ queries over the same corpus of $80$ tables.
Prior work has shown that increasing the number of retrieved candidates exhibits diminishing returns and introduces a trade-off between gold-table coverage and reasoning cost~\cite{craft}.
To quantify the incremental retrieval benefit, we define the marginal gain between two candidate sizes $k_1$ and $k_2$ as

\[
\operatorname{MG}(k_1,k_2)
=
\frac{\operatorname{Recall}@k_2-\operatorname{Recall}@k_1}
{k_2-k_1}.
\]

Both our retriever and CRAFT's retriever exhibit diminishing marginal gains as the candidate size increases, as shown in Table~\ref{tab:candidate_size_selection}.
Our method and OpenTab+r share the same initial hybrid retriever and therefore use the same retrieval curve for candidate-size selection.
For our retriever, Recall@$k$ increases from $0.61$ at $k=1$ to $0.97$ at $k=10$, while the marginal gain decreases from $0.1203$ between $k=1$ and $k=3$ to $0.0016$ between $k=10$ and $k=20$.
CRAFT exhibits a similar saturation pattern, with Recall@$k$ increasing from $0.55$ at $k=1$ to $0.91$ at $k=10$, followed by a substantially smaller marginal gain beyond $k=10$.

Based on the diminishing retrieval gains beyond $k=10$, we select $k=10$ as the operating point for our method and CRAFT.
We use $k=10$ for both OpenTab and OpenTab+r to ensure a consistent comparison between the two variants.

\section{Main Results}
\begin{table*}[t]
\centering
\caption{End-to-end answer accuracy on DataBench across different LLM backbones.
Bold indicates the best OTQA performance for each LLM configuration.
The closed-domain setting is reported as a diagnostic reference in which the gold table is directly provided to the reasoner.
A dash indicates that the corresponding experiment was not conducted.}
\label{tab:main_results}
\begin{tabular}{lccccc}
\toprule
\textbf{LLM} &
\textbf{Closed-domain} &
\textbf{Ours} &
\textbf{OpenTab} &
\textbf{OpenTab+r} &
\textbf{CRAFT} \\
\midrule
GPT-5.6-Sol
& \textit{0.86}
& \textbf{0.60}
& 0.17
& 0.36
& 0.27 \\

Gemini-3.7-Flash
& \textit{0.80}
& \textbf{0.55}
& --
& 0.37
& 0.28 \\

Qwen3.5-9B
& \textit{0.76}
& \textbf{0.61}
& --
& 0.24
& 0.20 \\

Qwen3.5-4B
& \textit{0.75}
& \textbf{0.59}
& --
& 0.22
& 0.22 \\
\bottomrule
\end{tabular}
\end{table*}


\paragraph{Our method consistently achieves the strongest end-to-end performance across all evaluated LLM backbones.}
As shown in Table~\ref{tab:main_results}, our method achieves the highest end-to-end accuracy among the evaluated OTQA methods across all four LLM backbones.
For GPT-5.6-Sol, our method achieves an accuracy of $0.60$, compared with $0.36$ for the strongest baseline, OpenTab+r.
Similar performance advantages are observed with Gemini-3.7-Flash and both Qwen3.5 models.
These results show that the performance advantage of our pipeline is consistent across different model families and scales.

\paragraph{Improving table retrieval alone is insufficient for strong OLTQA performance.}
Replacing OpenTab's original retriever with our table retriever improves its end-to-end accuracy from $0.17$ to $0.36$ under GPT-5.6-Sol.
This substantial improvement demonstrates the importance of table retrieval for end-to-end OLTQA performance.
However, OpenTab+r still remains considerably below our method at $0.60$, suggesting that retrieval improvement alone does not account for the remaining performance gap.

One limitation of OpenTab is that, after reranking the retrieved candidates, only the top-ranked table is retained for the reasoning stage.
Consequently, the reasoner cannot recover from a reranking error by considering other retrieved candidate tables.
In addition, OpenTab relies on a fixed table context during reasoning.
If query-relevant information is omitted from this context, the reasoner has no mechanism to dynamically access additional table information.
This limitation is relevant to large-table reasoning because a compact fixed context may omit information required to answer the query.

CRAFT provides further evidence that strong table retrieval alone does not guarantee strong end-to-end OLTQA performance.
As shown in Table~\ref{tab:retrieval_results}, CRAFT achieves retrieval recall comparable to our method, with Recall@10 of $0.92$ and $0.91$, respectively.
However, its end-to-end accuracy remains substantially lower across all evaluated LLM backbones (Table~\ref{tab:main_results}).
CRAFT also constructs a fixed table context for reasoning using all column names and the top-$5$ query-relevant rows from each candidate table.
Similar to OpenTab, once this context has been constructed, the reasoner cannot dynamically access table information omitted from the context.
In contrast, our reasoner can interactively access additional information from the candidate tables during reasoning.

Taken together, these results suggest that retrieving the gold table alone is insufficient for strong OLTQA performance.
Retaining multiple candidate tables during reasoning can improve gold-table coverage.
Moreover, even when the gold table is available to the reasoner, the provided table context must contain sufficient query-relevant information.
An insufficient table context may lack information required to answer the query and consequently lead to reasoning errors.
When the initially constructed table context is incomplete, providing a mechanism to access additional table information during reasoning allows the reasoner to obtain missing information as needed.

\section{Table Retrieval Results}

\begin{table*}[t]
\centering
\caption{Table retrieval performance on DataBench.
R@$k$ denotes the proportion of queries for which the gold table appears within the top-$k$ retrieved candidates.
Hybrid retrieval combines BM25 sparse retrieval with Transformer-based dense retrieval.
CRAFT uses a cascaded retrieval pipeline with stage-dependent retrieval strategies and table representations.}
\label{tab:retrieval_results}
\begin{tabular}{lllcccc}
\toprule
\textbf{Method} &
\textbf{Retriever} &
\textbf{Table Representation} &
\textbf{R@1} &
\textbf{R@3} &
\textbf{R@5} &
\textbf{R@10} \\
\midrule
OpenTab~\cite{opentab}
& BM25
& Full-content
& 0.04 & 0.07 & 0.09 & 0.14 \\

OpenTab (Hybrid)
& Hybrid
& Full-content
& 0.19 & 0.34 & 0.45 & 0.61 \\

CRAFT~\cite{craft}
& Cascaded
& Stage-dependent
& 0.56 & \textbf{0.78} & \textbf{0.86} & \textbf{0.92} \\
\midrule
Ours / OpenTab+r
& Hybrid
& Schema-level
& \textbf{0.65} & 0.77 & 0.83 & 0.91 \\
\bottomrule
\end{tabular}
\end{table*}

\paragraph{Both the retrieval algorithm and table representation substantially affect table retrieval performance.}
Table~\ref{tab:retrieval_results} compares the retrieval performance of different configurations on DataBench.
OpenTab's original BM25 retriever with a full-content table representation achieves only $0.04$ Recall@1 and $0.14$ Recall@10.
Replacing BM25 with a hybrid retrieval strategy while retaining the full-content table representation substantially improves Recall@1 to $0.19$ and Recall@10 to $0.61$, showing that the retrieval algorithm itself is an important factor.

However, replacing the full-content table representation with the schema-level table representation further increases Recall@1 to $0.65$ and Recall@10 to $0.91$.
This result suggests that including large amounts of cell-level content in the retrieval representation may introduce query-irrelevant information and weaken the table-level semantic signal, a concern that becomes particularly relevant for large tables due to their substantially larger volume of cell-level content.
In contrast, a schema-level table representation provides a more focused representation for table identification.

Overall, the results highlight the importance of both the retrieval algorithm and table representation for table retrieval in OLTQA.